\begin{proof}
Use the observable companion coordinates
\[
 T=\begin{pmatrix}0&1\\-1/2&-1\end{pmatrix},\quad
 l_0=(1,0),\quad
 Q=\begin{pmatrix}8/5&4/5\\4/5&12/5\end{pmatrix},
 \qquad Q-T^*QT=l_0^*l_0.
\]
The canonical Euclidean realization is
$D=Q^{1/2}TQ^{-1/2}$, $l=l_0Q^{-1/2}$.  Equivalently, compute in the
$Q$ inner product, in which
\[
 \Gamma=\frac{\sqrt7}{4},\qquad
 x_{\rm raw}=\frac{(-1,-3)^{\mathsf T}}{2\sqrt7},\quad
 a_0=-\frac1{2\sqrt7},\quad a_1=-\frac3{2\sqrt7},\quad a_2=\Gamma.
\]
Direct rational algebra gives
\[
 S=\frac{23\sqrt7}{196},\qquad
 v_{\rm raw}=\frac{(1,3)^{\mathsf T}}{28},\qquad
 K_{\rm raw}=\begin{pmatrix}
 2\sqrt7/7&-\sqrt7/7\\-\sqrt7/28&13\sqrt7/28
 \end{pmatrix}.
\]
The last matrix is $Q$-self-adjoint and represents the same positive
response operator.  In particular,
\[
 v_{\rm raw}^*QK_{\rm raw}^{-1}v_{\rm raw}=\frac{3\sqrt7}{196}.
\]
The three terms in \eqref{I:R:eq:bk-response-formula} are therefore
\[
 \frac9{28}-\frac{23}{112}+\frac3{112}=\frac17.
\]
The value for this extremizer follows without any uniqueness assertion.
The last statement follows by taking the maximum in
\eqref{I:R:eq:bk-response-formula}.
\end{proof}
